\bookchapter{Bab 63: View, Negeri (State) dan Navigasi}

\chapterintro{Mainan satu halaman sudah pandai dibuat. Bab ini menyelesaikan dua soalan 'membesar': cara melompat-lompat antara berbilang halaman (navigasi), dan cara data mengalir antara halaman (pengurusan state). Selepas bab ini, anda boleh membina kerangka App berbilang halaman klasik seperti 'senarai butiran sunting'.}

\booksection{Peta menyeluruh aliran data pengurusan state}

Alat state SwiftUI bukan setakat @State. Lihat dahulu peta menyeluruh ini untuk memahami 'bila guna yang mana':

\begin{center}
\begin{tabularx}{\linewidth}{|X|X|X|X|}
\hline
\textbf{Alat} & \textbf{Mengurus apa} & \textbf{Digunakan oleh} & \textbf{Senario tipikal} \tabularnewline
\hline
@State & State sementara peribadi halaman & View semasa & Pengira, teks medan input \tabularnewline
@Binding & Meminjam state orang lain & Sub-view mengubah view induk & Komponen tersuai \tabularnewline
@StateObject / @Observable & Logik dikongsi merentas halaman & Model pandangan (ViewModel) & Pemuatan data, logik perniagaan \tabularnewline
@EnvironmentObject / .environment & State dikongsi secara global & Keseluruhan App & Status log masuk pengguna, tetapan tema \tabularnewline
\hline
\end{tabularx}
\end{center}

\begin{center}\small\color{gray}Peta menyeluruh alat state SwiftUI (pada 2026, @Observable sudah arus perdana)\end{center}

\begin{lstlisting}[language=Swift]
// Cara moden iOS 17+: makro @Observable (digalakkan untuk projek baru)
import Observation

@Observable
class CounterModel {                 // View model
    var count = 0
    func increment() { count += 1 }
}

struct CounterView: View {
    @State private var model = CounterModel()   // Memegang model

    var body: some View {
        VStack {
            Text("Anda menekan \(model.count) kali")
            Button("Tekan saya") { model.increment() }
        }
    }
}
\end{lstlisting}

Nasihat terus untuk pemula: state peribadi halaman guna @State, yang merentas halaman atau berlogik guna kelas @Observable. @StateObject/@ObservableObject dalam tutorial lama ialah gaya generasi sebelumnya — cukup sekadar boleh membaca; untuk kod baharu, ikut seperti di atas.

\booksection{NavigationStack: 'tindanan' lompatan halaman}

Mulai iOS 16, navigasi berbilang halaman menggunakan NavigationStack — seperti timbunan pinggan: satu dimasukkan satu demi satu, dan dikeluarkan satu demi satu semasa kembali:

\begin{lstlisting}[language=Swift]
struct HomeView: View {
    var body: some View {
        NavigationStack {
            List {
                NavigationLink("Buku akaun Zhang San") {
                    DetailView(name: "Buku akaun Zhang San")   // halaman destinasi
                }
                NavigationLink("Buku akaun Li Si") {
                    DetailView(name: "Buku akaun Li Si")
                }
            }
            .navigationTitle("Senarai buku akaun")   // tajuk besar di atas
        }
    }
}

struct DetailView: View {
    let name: String
    var body: some View {
        Text("Halaman butiran \(name)")
            .navigationTitle(name)
            .navigationBarTitleDisplayMode(.inline)
    }
}
\end{lstlisting}

NavigationLink ialah 'pautan lompatan yang boleh ditekan'; navigationTitle menjana tajuk bar navigasi dan butang kembali secara automatik. Untuk lompatan secara berprogram (contohnya kembali automatik selepas menekan 'Simpan'), gunakan @Environment(\textbackslash{}.dismiss) untuk menutup halaman semasa — dengan dua helah ini, 90\% keperluan navigasi sudah lengkap.

\booksection{Pop timbul dan tab bawah: sheet dan TabView}

Selain 'menolak halaman', ada dua lagi interaksi klasik:

\begin{lstlisting}[language=Swift]
// 1. Pop separuh skrin/skrin penuh dari bawah (kegemaran borang input)
struct HomeView: View {
    @State private var showAdd = false
    var body: some View {
        Button("Catat") { showAdd = true }
            .sheet(isPresented: $showAdd) {
                AddRecordView()     // Halaman input yang dipaparkan
            }
    }
}

// 2. Bar tab bawah ("rangka" aplikasi)
struct MainTabView: View {
    var body: some View {
        TabView {
            RecordListView()
                .tabItem { Label("Buku akaun", systemImage: "book") }
            StatsView()
                .tabItem { Label("Statistik", systemImage: "chart.pie") }
            SettingsView()
                .tabItem { Label("Tetapan", systemImage: "gear") }
        }
    }
}
\end{lstlisting}

\begin{tcbinfo}{Petua: Bila guna yang mana?}
Peraturan pengalaman: aliran utama guna NavigationStack untuk menolak halaman (senaraibutiran), operasi sementara guna pop timbul sheet (tambah, sunting, sahkan), beberapa bahagian besar yang setara guna TabView (tab bawah). Gabungan ketiga-tiganya ialah kerangka hampir semua App.

\end{tcbinfo}

\booksection{List: senarai paling popular di iOS}

Komponen paling kerap digunakan dan paling 'berbau Apple' di iOS ialah List — sedia ada garis pemisah, padam dengan leretan dan sorotan semasa ditekan:

\begin{lstlisting}[language=Swift]
struct RecordListView: View {
    @State private var records = ["Makan tengah hari 25.5", "LRT 6.0", "Gaji +8000"]

    var body: some View {
        List {
            ForEach(records, id: \.self) { r in
                Text(r)
            }
            .onDelete { indexSet in          // leret kiri untuk padam
                records.remove(atOffsets: indexSet)
            }
        }
    }
}
\end{lstlisting}

Senarai + NavigationLink + sheet + TabView — gabungkan empat elemen ini, dan anda sudah boleh membina kerangka App lengkap 'senarai rekod buka butiran tekan + muncul borang kemasukan'. Projek buku rekod perbelanjaan bab seterusnya dibina begini.

\begin{tcbpractice}{Praktik: Bina 'kerangka tiga halaman'}
\booksection{Kitaran hayat aplikasi: hadapan dan latar belakang}

Sama seperti kitaran hayat Activity pada Android, App iOS juga mengalami perubahan state 'lahir, hadapan, latar belakang, mati'. Dalam SwiftUI, gunakan nilai persekitaran scenePhase untuk memerhatikannya:

\begin{lstlisting}[language=Swift]
@main
struct MyApp: App {
    @Environment(\.scenePhase) private var phase

    var body: some Scene {
        WindowGroup {
            ContentView()
        }
        .onChange(of: phase) { _, newPhase in
            switch newPhase {
            case .active:    print("Kembali ke hadapan, sambung segar semula")
            case .inactive:  print("Dalam peralihan (panggilan masuk/baru beralih ke latar)")
            case .background: print("Masuk latar, simpan draf")
            @unknown default: break
            }
        }
    }
}
\end{lstlisting}

Kegunaan tipikal: segarkan data semasa kembali ke hadapan, simpan draf semasa masuk ke latar belakang. iOS mengawal pelaksanaan latar belakang dengan sangat ketat — App biasa akan segera digantung (dibekukan) selepas masuk latar belakang; jangan berharap boleh membuat kerja berat secara senyap-senyap di latar belakang. Jika benar-benar perlu tugasan berkala, gunakan pemberitahuan tempatan dalam Bab 68, bukan 'kekal berjalan di latar belakang'.

Bina kerangka versi data palsu buku rekod perbelanjaan: TabView tiga tab (Rekod/Statistik/Tetapan) halaman rekod guna List memaparkan data palsu, tekan baris untuk ke butiran (NavigationStack) butang '+' penjuru atas kanan memunculkan sheet kemasukan. Tanpa menyimpan data, tanpa internet — semata-mata latihan kerangka. Kerangka ini siap, mengisinya dengan data sebenar pada Bab 65 menjadi perkara semula jadi.

\end{tcbpractice}

\begin{tcbsummary}{Ringkasan bab ini}
\begin{itemize}
  \item Empat alat state: @State (peribadi halaman), @Binding (pinjaman), @Observable (logik), .environment (global)
  \item Projek baharu guna makro @Observable menulis ViewModel; @ObservableObject tutorial lama cukup sekadar dapat dibaca
  \item NavigationStack menolak halaman; @Environment(\textbackslash{}.dismiss) kembali secara berprogram
  \item Aliran utama tolak halaman, operasi sementara munculkan sheet, bahagian besar setara guna TabView
  \item List ialah senarai paling popular: garis pemisah dan padam leretan sedia guna
\end{itemize}

\end{tcbsummary}
